%% ----------------------------------------------------------------------------
%%  §2.5  Il binomio di Newton
%% ----------------------------------------------------------------------------
\section{Il binomio di Newton}
\label{sec:binomio}

Un'importante applicazione dei risultati del paragrafo precedente (che peraltro è
anche conseguenza del principio di induzione, come mostreremo nella seconda parte del
paragrafo) è la formula che esprime la potenza $n$-esima di un binomio. Per scriverla in
modo compatto conviene introdurre un simbolo che useremo in tutto il resto del libro.

\begin{notazione}[Il simbolo di sommatoria]\index{sommatoria}
Dati i numeri reali $c_{p},c_{p+1},\dots,c_{q}$ (con $p\le q$ interi), la loro somma si
indica con
\[
  \sum_{k=p}^{q}c_{k}=c_{p}+c_{p+1}+\dots+c_{q},
\]
che si legge «somma per $k$ che va da $p$ a $q$ di $c_{k}$». La lettera $k$ è un
\emph{indice muto}: $\sum_{k=p}^{q}c_{k}$ e $\sum_{j=p}^{q}c_{j}$ indicano la stessa
somma. Valgono le proprietà (che seguono dalle proprietà associativa, commutativa e
distributiva)
\[
  \sum_{k=p}^{q}(c_{k}+d_{k})=\sum_{k=p}^{q}c_{k}+\sum_{k=p}^{q}d_{k}, \qquad
  \sum_{k=p}^{q}h\,c_{k}=h\sum_{k=p}^{q}c_{k}\quad(h\in\R),
\]
e la regola di \emph{traslazione dell'indice}: ponendo $j=k+1$,
\begin{equation}\label{eq:traslazione-indice}
  \sum_{k=p}^{q}c_{k}=\sum_{j=p+1}^{q+1}c_{j-1}.
\end{equation}
Con questa notazione, ad esempio, la \eqref{eq:gauss} si scrive
$\sum_{k=1}^{n}k=\frac{n(n+1)}{2}$.
\end{notazione}

\begin{teorema}[Formula del binomio di Newton]\label{teo:binomio}\index{binomio di Newton}
Per ogni coppia di numeri reali $a,b$ e per ogni $n\in\N$ vale l'identità
\begin{equation}\label{eq:binomio}
\begin{split}
  (a+b)^{n}&=\binom n0a^{n}+\binom n1a^{n-1}b+\binom n2a^{n-2}b^{2}+\dots\\
  &\qquad\dots+\binom nk a^{n-k}b^{k}+\dots+\binom{n}{n-1}ab^{n-1}+\binom nn b^{n}\\
  &=\sum_{k=0}^{n}\binom nk a^{n-k}b^{k}.
\end{split}
\end{equation}
\end{teorema}

\begin{schema}
\Ipotesi{$a,b\in\R$; $n\in\N$ (con la convenzione $c^{0}=1$ per ogni $c$, anche
$c=0$, nei termini estremi della somma).}
\Tesi{$(a+b)^{n}=\displaystyle\sum_{k=0}^{n}\binom nk a^{n-k}b^{k}$.}
\end{schema}

Diamo due dimostrazioni, che corrispondono ai due modi di pensare descritti nel
paragrafo precedente: \emph{contare} e \emph{calcolare}.

\begin{strategia}[Piano d'attacco della prima dimostrazione]
Scriviamo $(a+b)^{n}$ come prodotto di $n$ fattori uguali e immaginiamo di svolgerlo
con la proprietà distributiva senza raccogliere nulla: otteniamo una somma di $2^{n}$
prodotti, ciascuno dei quali si forma scegliendo da ogni parentesi o la $a$ o la $b$
(\cref{fig:binomio-scelte}). Un prodotto in cui la $b$ è stata scelta $k$ volte vale
$a^{n-k}b^{k}$. Quanti sono? Tanti quanti i modi di scegliere \emph{quali} $k$ parentesi
forniscono la $b$, cioè i sottoinsiemi di $k$ elementi di un insieme di $n$ parentesi:
$\binom nk$. Il coefficiente binomiale deve il suo nome proprio a questo.
\end{strategia}

\begin{dimostrazione}[Prima dimostrazione (combinatoria)]
Immaginiamo di eseguire il prodotto di $n$ fattori
\begin{equation}\label{eq:binomio-prodotto}
  (a+b)^{n}=(a+b)\cdot(a+b)\cdots(a+b) \qquad (n\ \text{fattori});
\end{equation}
applicando ripetutamente la proprietà distributiva, il risultato si scrive come somma
algebrica di addendi, ciascuno dei quali si ottiene scegliendo un termine ($a$ oppure
$b$) da ognuno degli $n$ fattori e moltiplicando i termini scelti. Fissiamo la nostra
attenzione sugli addendi del tipo $a^{n-k}b^{k}$: eseguendo la moltiplicazione indicata,
si ottiene tale addendo tante volte quante è possibile scegliere i $k$ fattori dai quali
prendere la $b$ tra gli $n$ dati (dagli altri $n-k$ si prende necessariamente la $a$).
Per il \cref{teo:combinazioni}, ciò avviene $\binom nk$ volte. Raccogliendo gli addendi
uguali si ottiene la \eqref{eq:binomio}.
\end{dimostrazione}

\begin{marginfigure}
\centering
\small
\renewcommand{\arraystretch}{1.15}
\begin{tabular}{@{}c@{\quad}l@{\quad}c@{}}
\toprule
n.\ di $b$ & prodotti & tot.\\
\midrule
$0$ & $aaa$ & $1$\\
$1$ & $aab,\ aba,\ baa$ & $3$\\
$2$ & $abb,\ bab,\ bba$ & $3$\\
$3$ & $bbb$ & $1$\\
\bottomrule
\end{tabular}
\caption{Gli $2^{3}=8$ prodotti che si ottengono svolgendo $(a+b)(a+b)(a+b)$, raggruppati
secondo il numero di $b$ scelte: $(a+b)^{3}=a^{3}+3a^{2}b+3ab^{2}+b^{3}$.}
\label{fig:binomio-scelte}
\end{marginfigure}

La formula del binomio \eqref{eq:binomio} fornisce in particolare risultati di facile
verifica per $n=2$ e $n=3$: infatti, per tali valori di $n$ si ottengono le ben note
relazioni
\begin{gather}
  \binom20=1,\quad \binom21=2,\quad \binom22=1, \notag\\
  \text{da cui}\quad (a+b)^{2}=a^{2}+2ab+b^{2};\\[3pt]
  \binom30=\binom33=1,\quad \binom31=\binom32=3, \notag\\
  \text{da cui}\quad (a+b)^{3}=a^{3}+3a^{2}b+3ab^{2}+b^{3}.
\end{gather}

\subsection{La formula di Stifel e il triangolo di Tartaglia}

Le considerazioni che seguono sono utili per sviluppare esplicitamente il secondo
membro della formula del binomio per valori di $n$ più grandi di $3$, e forniscono
l'ingrediente della seconda dimostrazione.

\begin{lemma}[Formula di Stifel]\label{lem:stifel}\index{formula di Stifel}
Per ogni $n\in\N$ e ogni $k\in\N$ con $1\le k\le n$ vale la formula
\begin{equation}\label{eq:stifel}
  \binom{n}{k-1}+\binom nk=\binom{n+1}{k}.
\end{equation}
\end{lemma}

\begin{schema}
\Ipotesi{$n\in\N$, \ $1\le k\le n$.}
\Tesi{$\binom{n}{k-1}+\binom nk=\binom{n+1}{k}$.}
\end{schema}

\begin{strategia}
È un calcolo con i fattoriali: si mettono le due frazioni a denominatore comune
raccogliendo i fattori comuni $n!$, $(k-1)!$ e $(n-k)!$, e si usano le relazioni
$k!=k\cdot(k-1)!$ e $(n-k+1)!=(n-k+1)\cdot(n-k)!$. Dopo il calcolo daremo anche
un'interpretazione combinatoria, che spiega l'identità senza conti.
\end{strategia}

\begin{dimostrazione}
Dalla definizione \eqref{eq:combinazioni} otteniamo
\begin{equation}\label{eq:stifel-calcolo}
\begin{split}
  \binom{n}{k-1}+\binom nk
  &=\frac{n!}{(k-1)!\,(n-k+1)!}+\frac{n!}{k!\,(n-k)!}\\
  &=\frac{n!}{(k-1)!\,(n-k)!}\Bigl(\frac{1}{n-k+1}+\frac1k\Bigr)\\
  &=\frac{n!}{(k-1)!\,(n-k)!}\cdot\frac{n+1}{k\,(n-k+1)}\\
  &=\frac{(n+1)!}{k!\,(n+1-k)!}=\binom{n+1}{k}.
\end{split}
\end{equation}
Il calcolo è lecito anche per $k=n$, grazie alla convenzione $0!=1$.
\end{dimostrazione}

\begin{osservazione}[Interpretazione combinatoria della formula di Stifel]
Consideriamo un insieme di $n+1$ elementi $\{a_{1},\dots,a_{n},a_{n+1}\}$ e contiamo i
suoi sottoinsiemi di $k$ elementi, che sono $\binom{n+1}{k}$, dividendoli in due gruppi:
quelli che \emph{non} contengono $a_{n+1}$, che sono i sottoinsiemi di $k$ elementi di
$\{a_{1},\dots,a_{n}\}$, cioè $\binom nk$; e quelli che \emph{contengono} $a_{n+1}$, che
si ottengono aggiungendo $a_{n+1}$ a un sottoinsieme di $k-1$ elementi di
$\{a_{1},\dots,a_{n}\}$, cioè $\binom{n}{k-1}$. La somma è la \eqref{eq:stifel}.
\end{osservazione}

La proprietà \eqref{eq:stifel} permette di scrivere il seguente \emph{triangolo di
Tartaglia} (noto anche come \emph{triangolo di Pascal}), nel quale ogni coefficiente è
uguale alla somma dei due coefficienti più vicini della riga precedente
(\cref{fig:tartaglia}).

\begin{figure}[htb]
\centering
\begin{tikzpicture}[x=0.82cm, y=0.62cm, every node/.style={font=\footnotesize\addfontfeatures{Numbers=Lining}},
  cella/.style={circle, inner sep=1.2pt, minimum size=5.4mm, fill=bluchiaro, draw=blu!40},
  celrame/.style={cella, fill=rame!30, draw=rame},
  celverde/.style={cella, fill=verde!30, draw=verde}]
  % (valori calcolati una volta per tutte: pgfmath non gestisce i fattoriali grandi)
  \node[grigio, anchor=east] at (-5.2,-0) {$n=0$};
  \node[cella] (t-0-0) at (0,-0) {1};
  \node[grigio, anchor=east] at (-5.2,-1) {$n=1$};
  \node[cella] (t-1-0) at (-0.5,-1) {1};
  \node[cella] (t-1-1) at (0.5,-1) {1};
  \node[grigio, anchor=east] at (-5.2,-2) {$n=2$};
  \node[cella] (t-2-0) at (-1,-2) {1};
  \node[cella] (t-2-1) at (0,-2) {2};
  \node[cella] (t-2-2) at (1,-2) {1};
  \node[grigio, anchor=east] at (-5.2,-3) {$n=3$};
  \node[cella] (t-3-0) at (-1.5,-3) {1};
  \node[cella] (t-3-1) at (-0.5,-3) {3};
  \node[cella] (t-3-2) at (0.5,-3) {3};
  \node[cella] (t-3-3) at (1.5,-3) {1};
  \node[grigio, anchor=east] at (-5.2,-4) {$n=4$};
  \node[cella] (t-4-0) at (-2,-4) {1};
  \node[cella] (t-4-1) at (-1,-4) {4};
  \node[cella] (t-4-2) at (0,-4) {6};
  \node[cella] (t-4-3) at (1,-4) {4};
  \node[cella] (t-4-4) at (2,-4) {1};
  \node[grigio, anchor=east] at (-5.2,-5) {$n=5$};
  \node[cella] (t-5-0) at (-2.5,-5) {1};
  \node[cella] (t-5-1) at (-1.5,-5) {5};
  \node[celrame] (t-5-2) at (-0.5,-5) {10};
  \node[celrame] (t-5-3) at (0.5,-5) {10};
  \node[cella] (t-5-4) at (1.5,-5) {5};
  \node[cella] (t-5-5) at (2.5,-5) {1};
  \node[grigio, anchor=east] at (-5.2,-6) {$n=6$};
  \node[cella] (t-6-0) at (-3,-6) {1};
  \node[cella] (t-6-1) at (-2,-6) {6};
  \node[cella] (t-6-2) at (-1,-6) {15};
  \node[celverde] (t-6-3) at (0,-6) {20};
  \node[cella] (t-6-4) at (1,-6) {15};
  \node[cella] (t-6-5) at (2,-6) {6};
  \node[cella] (t-6-6) at (3,-6) {1};
  \node[grigio, anchor=east] at (-5.2,-7) {$n=7$};
  \node[cella] (t-7-0) at (-3.5,-7) {1};
  \node[cella] (t-7-1) at (-2.5,-7) {7};
  \node[celrame] (t-7-2) at (-1.5,-7) {21};
  \node[celrame] (t-7-3) at (-0.5,-7) {35};
  \node[cella] (t-7-4) at (0.5,-7) {35};
  \node[cella] (t-7-5) at (1.5,-7) {21};
  \node[cella] (t-7-6) at (2.5,-7) {7};
  \node[cella] (t-7-7) at (3.5,-7) {1};
  \node[grigio, anchor=east] at (-5.2,-8) {$n=8$};
  \node[cella] (t-8-0) at (-4,-8) {1};
  \node[cella] (t-8-1) at (-3,-8) {8};
  \node[cella] (t-8-2) at (-2,-8) {28};
  \node[celverde] (t-8-3) at (-1,-8) {56};
  \node[cella] (t-8-4) at (0,-8) {70};
  \node[cella] (t-8-5) at (1,-8) {56};
  \node[cella] (t-8-6) at (2,-8) {28};
  \node[cella] (t-8-7) at (3,-8) {8};
  \node[cella] (t-8-8) at (4,-8) {1};
  \draw[-{Latex[length=1.6mm]}, rame, thick] (t-5-2) -- (t-6-3);
  \draw[-{Latex[length=1.6mm]}, rame, thick] (t-5-3) -- (t-6-3);
  \draw[-{Latex[length=1.6mm]}, rame, thick] (t-7-2) -- (t-8-3);
  \draw[-{Latex[length=1.6mm]}, rame, thick] (t-7-3) -- (t-8-3);
  \node[rame, anchor=west] at (5.0,-5.5) {$\binom52+\binom53=\binom63$};
  \node[rame, anchor=west] at (5.0,-7.5) {$\binom72+\binom73=\binom83$};
\end{tikzpicture}
\caption{Il triangolo di Tartaglia: nella riga $n$ compaiono i coefficienti
$\binom n0,\binom n1,\dots,\binom nn$. Ogni numero è la somma dei due che gli stanno
sopra (formula di Stifel): ad esempio $10+10=20$ e $21+35=56$. La simmetria di ogni riga
è la \eqref{eq:simmetria-binomiale}.}
\label{fig:tartaglia}
\end{figure}

Quindi, ad esempio, leggendo la riga $n=5$ risulta
\begin{equation}\label{eq:binomio-5}
  (a+b)^{5}=a^{5}+5a^{4}b+10a^{3}b^{2}+10a^{2}b^{3}+5ab^{4}+b^{5}.
\end{equation}

\subsection{La dimostrazione per induzione}

Di seguito esponiamo un'ulteriore dimostrazione della formula del binomio, con il metodo
di induzione.

\needspace{8\baselineskip}
\begin{strategia}[Piano d'attacco della seconda dimostrazione]
Supposta vera la formula per $n$, otteniamo $(a+b)^{n+1}$ moltiplicando lo sviluppo di
$(a+b)^{n}$ per $(a+b)$. Si ottengono due somme: in quella che proviene dalla
moltiplicazione per $a$ l'esponente di $a$ sale di uno, in quella che proviene dalla
moltiplicazione per $b$ sale l'esponente di $b$. Per poterle sommare termine a termine
occorre che i termini con la stessa potenza $a^{n+1-k}b^{k}$ abbiano lo stesso indice:
lo si ottiene traslando l'indice della seconda somma \eqref{eq:traslazione-indice}. A
quel punto ogni potenza $a^{n+1-k}b^{k}$ ha coefficiente $\binom nk+\binom n{k-1}$, che
per la formula di Stifel è proprio $\binom{n+1}k$. È, in formule, la costruzione del
triangolo di Tartaglia: la riga $n+1$ si ottiene dalla riga $n$.
\end{strategia}

\begin{dimostrazione}[Seconda dimostrazione (per induzione)]
Cominciamo con l'osservare che la formula del binomio è verificata per $n=1$; infatti
\begin{equation}\label{eq:binomio-base}
  (a+b)^{1}=\binom10a+\binom11b=a+b.
\end{equation}
Supponiamo vera la \eqref{eq:binomio} per un indice $n\in\N$ e dimostriamo la validità
della formula analoga con $n+1$ al posto di $n$. A tal fine moltiplichiamo entrambi i
membri della \eqref{eq:binomio} per $(a+b)$, ottenendo
\[
  (a+b)^{n+1}=\sum_{k=0}^{n}\binom nk a^{n+1-k}b^{k}+\sum_{k=0}^{n}\binom nk a^{n-k}b^{k+1}.
\]
Nella seconda somma poniamo $j=k+1$ (\eqref{eq:traslazione-indice}); poiché
$a^{n-k}b^{k+1}=a^{n+1-j}b^{j}$, essa diventa
$\sum_{j=1}^{n+1}\binom{n}{j-1}a^{n+1-j}b^{j}$. Rinominando $j$ in $k$ (l'indice è muto)
e isolando il primo termine della prima somma ($k=0$) e l'ultimo della seconda
($k=n+1$):
\begin{equation}\label{eq:binomio-passo}
\begin{split}
  (a+b)^{n+1}&=a^{n+1}+\sum_{k=1}^{n}\Bigl[\binom nk+\binom n{k-1}\Bigr]a^{n+1-k}b^{k}+b^{n+1}\\
  &=\binom{n+1}{0}a^{n+1}+\sum_{k=1}^{n}\binom{n+1}{k}a^{n+1-k}b^{k}+\binom{n+1}{n+1}b^{n+1}\\
  &=\sum_{k=0}^{n+1}\binom{n+1}{k}a^{n+1-k}b^{k},
\end{split}
\end{equation}
dove abbiamo usato la formula di Stifel \eqref{eq:stifel} per $1\le k\le n$ e la
\eqref{eq:binomiale-estremi}. È la \eqref{eq:binomio} con $n+1$ al posto di $n$; per il
principio di induzione la formula è dimostrata.
\end{dimostrazione}

\subsection{Conseguenze e applicazioni}

\begin{corollario}\label{cor:binomio}
Per ogni $n\in\N$:
\begin{enumerate}
  \item $\displaystyle\sum_{k=0}^{n}\binom nk=2^{n}$;
  \item $\displaystyle\sum_{k=0}^{n}(-1)^{k}\binom nk=0$;
  \item se $x\ge0$, \ $(1+x)^{n}\ge1+nx+\dfrac{n(n-1)}{2}\,x^{2}$.
\end{enumerate}
\end{corollario}

\begin{schema}
\Ipotesi{$n\in\N$; \ per la (iii), $x\ge0$.}
\Tesi{(i) $\sum_{k=0}^{n}\binom nk=2^{n}$; \ (ii) $\sum_{k=0}^{n}(-1)^{k}\binom nk=0$; \
(iii) $(1+x)^{n}\ge1+nx+\frac{n(n-1)}{2}x^{2}$.}
\end{schema}

\begin{strategia}
La formula del binomio è un'identità in $a$ e $b$: scegliendo valori particolari si
ottengono informazioni sui coefficienti. Con $a=b=1$ le potenze scompaiono e restano i
soli coefficienti; con $b=-1$ compaiono i segni alterni. Per la (iii) si sfrutta il
fatto che, se $b=x\ge0$, tutti gli addendi sono non negativi, e una somma di termini non
negativi non è minore di una parte di essi.
\end{strategia}

\begin{dimostrazione}
(i) e (ii) si ottengono dalla \eqref{eq:binomio} con $a=b=1$ e con $a=1$, $b=-1$. Per
(iii), con $a=1$ e $b=x\ge0$ tutti gli addendi della \eqref{eq:binomio} sono non
negativi; trascurando quelli con $k\ge3$ resta
$(1+x)^{n}\ge1+\binom n1x+\binom n2x^{2}=1+nx+\frac{n(n-1)}2x^{2}$ (per $n=1$ il
termine in $x^{2}$ è nullo e la disuguaglianza è un'uguaglianza).
\end{dimostrazione}

\begin{osservazione}
La (i) ha un significato combinatorio: $\binom nk$ conta i sottoinsiemi con $k$
elementi di un insieme di $n$ elementi, e sommando su tutti i $k$ si contano tutti i
sottoinsiemi, che sono $2^{n}$ (esercizio del \cref{sec:induzione}). La (iii) è un
rafforzamento della disuguaglianza di Bernoulli (\cref{teo:bernoulli}) per $x\ge0$, e
ci servirà nello studio delle successioni: per esempio, per dimostrare che
$\sqrt[n]{n}$ si avvicina a $1$ al crescere di $n$.
\end{osservazione}

\begin{esempio}[Sensibilità della portanza e della potenza alla velocità]\label{es:sensibilita}
La portanza di un'ala è $L=\frac12\rho V^{2}SC_{L}$. Se la velocità aumenta del $2\%$,
cioè passa da $V$ a $V(1+\epsilon)$ con $\epsilon=0{,}02$, a parità di $C_{L}$ la
portanza diventa
\[
  L(1+\epsilon)^{2}=L\,(1+2\epsilon+\epsilon^{2})=L\,(1+0{,}04+0{,}0004),
\]
cioè aumenta del $4{,}04\%$. La potenza necessaria a vincere la resistenza è
$P=DV=\frac12\rho V^{3}SC_{D}$; un aumento del $3\%$ della velocità ($\epsilon=0{,}03$)
la moltiplica per
\[
  (1+\epsilon)^{3}=1+3\epsilon+3\epsilon^{2}+\epsilon^{3}=1+0{,}09+0{,}0027+0{,}000027
  =1{,}092727,
\]
cioè più del $9\%$. La formula del binomio giustifica la regola pratica, usata
continuamente nelle stime ingegneristiche, $(1+\epsilon)^{n}\approx1+n\epsilon$ per
$\epsilon$ «piccolo»: l'errore commesso è costituito dai termini con $\epsilon^{2}$,
$\epsilon^{3}$, \dots, che per $\epsilon=0{,}03$ valgono complessivamente meno di
$0{,}003$. Una variazione percentuale piccola su una grandezza elevata alla $n$ produce,
in prima approssimazione, una variazione percentuale $n$ volte maggiore.
\end{esempio}

\begin{notastorica}[Tartaglia, Pascal, Newton --- e prima di loro]
Il triangolo porta in Italia il nome di Niccolò Fontana detto Tartaglia (Brescia,
1499 circa -- Venezia, 1557), che lo pubblicò nel \emph{General trattato di numeri et misure}
(1556), e in Francia quello di Blaise Pascal, che ne fece uno studio sistematico nel
\emph{Traité du triangle arithmétique} (scritto nel 1654, pubblicato postumo nel 1665),
usandolo per risolvere problemi di gioco
d'azzardo posti dal cavaliere de Méré: è l'atto di nascita del calcolo delle
probabilità. Ma il triangolo era noto da secoli: in Cina compare nel trattato di Yang Hui
(1261), che lo attribuisce a Jia Xian (XI secolo); in Persia era conosciuto da al-Karaji
e da Omar Khayyam; in India, in forma combinatoria, già nell'opera di Pingala sulla
metrica dei versi. La formula di Stifel prende il nome da Michael Stifel
(\emph{Arithmetica integra}, 1544). Il merito di Isaac Newton (1665) fu di estendere la
formula agli esponenti \emph{non interi}, ottenendo uno sviluppo in serie infinita di
$(1+x)^{\alpha}$ che ritroveremo studiando la formula di Taylor.
\end{notastorica}

\begin{python}[Il triangolo di Tartaglia costruito con la formula di Stifel]
import sympy as sp

def tartaglia(righe):
    riga = [1]
    for n in range(righe):
        yield riga
        # formula di Stifel: ogni elemento interno è la somma
        # dei due elementi che gli stanno sopra
        riga = ([1] + [riga[k - 1] + riga[k]
                       for k in range(1, len(riga))] + [1])

for riga in tartaglia(8):
    print(" ".join(f"{c:3d}" for c in riga).center(40))

a, b = sp.symbols("a b")
print(sp.expand((a + b)**5))
\end{python}
\begin{risultato}
                    1                   
                  1   1                 
                1   2   1               
              1   3   3   1             
            1   4   6   4   1           
          1   5  10  10   5   1         
        1   6  15  20  15   6   1       
      1   7  21  35  35  21   7   1     
a**5 + 5*a**4*b + 10*a**3*b**2 + 10*a**2*b**3 + 5*a*b**4 + b**5
\end{risultato}

\begin{figure}[htb]
\centering
\includegraphics[width=0.78\textwidth]{sierpinski}
\caption{Le prime $128$ righe del triangolo di Tartaglia, colorando i coefficienti
dispari e lasciando in bianco i pari (\texttt{python/fig\_sierpinski.py}). Compare il
\emph{triangolo di Sierpiński}, un frattale: la struttura si ripete identica a scale
diverse. Il motivo è proprio la formula di Stifel, letta «modulo $2$»: pari$+$pari e
dispari$+$dispari danno pari, pari$+$dispari dà dispari.}
\label{fig:sierpinski}
\end{figure}
